\documentclass[11pt]{article}

\usepackage[utf8]{inputenc}
\usepackage[T1]{fontenc}
\usepackage[margin=1in]{geometry}
\usepackage{graphicx}
\usepackage{booktabs}
\usepackage{longtable}
\usepackage{tabularx}
\usepackage{array}
\usepackage{enumitem}
\usepackage{xcolor}
% hypertexnames=false avoids duplicate destination warnings when page numbering/counters reset.
\usepackage[hypertexnames=false]{hyperref}
\usepackage{float}
\usepackage{tikz}
\usetikzlibrary{arrows.meta,fit,positioning,shapes.geometric}
\usepackage{listings}

\hypersetup{
  colorlinks=true,
  linkcolor=blue,
  urlcolor=blue,
  citecolor=blue
}

\setlength{\parindent}{0pt}
\setlength{\parskip}{0.35em}

\newcommand{\ProjectName}{HealthSync}
\newcommand{\CourseName}{S26CS6.401 Software Engineering}
\newcommand{\TeamNumber}{2}

% Visible TODO markers (intentionally rendered in the PDF so we can fill them before submission).
\newcommand{\TODO}[1]{\textcolor{red}{\textbf{TODO:} #1}}
% Visible "re-check / rephrase" marker.
\newcommand{\RECHECK}[1]{\textcolor{orange}{\textbf{RECHECK:} #1}}
% Diagram placeholder marker.
\newcommand{\DIAGRAM}[1]{\textcolor{violet}{\textbf{[DIAGRAM NEEDED]:} #1}}

\lstset{
  basicstyle=\ttfamily\small,
  columns=fullflexible,
  breaklines=true,
  frame=single,
  framerule=0.3pt,
  xleftmargin=0.6em,
  xrightmargin=0.6em
}

\begin{document}

% ===========================
% Title Page
% ===========================
\begin{titlepage}
  \centering
  {\Large \CourseName \par}
  \vspace{0.35cm}
  {\Large \textbf{Project 3 Report} \par}
  \vspace{1.0cm}
  {\huge \textbf{\ProjectName} \par}
  \vspace{0.25cm}
  {\Large Team \TeamNumber \par}
  \vspace{0.8cm}

  \begin{tabular}{@{}ll@{}}
    \toprule
    \textbf{Team Member} & \textbf{Roll Number} \\
    \midrule
    Kartik Vij          & 2022102060 \\
    Kushagra Dhingra    & 2022115004 \\
    Lakshmi Dedeepya    & 2024801007 \\
    Prithvi Karthik     & 2022101020 \\
    Samagra Bharti      & 2022115007 \\
    \bottomrule
  \end{tabular}

  \vfill
  \textbf{Repository link:} https://github.com/Prithvi-k/SE-Project-3 \par
  \vspace{0.35cm}
  \textbf{Date:} \today \par
\end{titlepage}

\pagenumbering{roman}
{\setlength{\parskip}{0pt}\tableofcontents}
\newpage
\pagenumbering{arabic}

% ===========================
% 1) Overview
% ===========================
\section{Overview}

\ProjectName\ is an integrated health tracking platform that combines meal logging, workout tracking,
health goal management, progress analytics, and personalized workout recommendations in a single system.
The motivation is to reduce the fragmentation users face when they rely on multiple apps to track diet,
exercise, and fitness goals.

% \subsection{Prototype focus for Project 3}
% For the Project 3 deliverables, we implement and document one primary end-to-end nontrivial flow:
% \textbf{authenticated meal logging with optional image-based meal analysis}.
% This flow exercises authentication, validation, API boundaries, persistence, UI/UX, and testing,
% and serves as the spine on which the rest of the architecture is demonstrated.

% \subsection{Tech stack}
% \begin{itemize}[leftmargin=1.2em]
%   \item \textbf{Frontend:} React 18 + TypeScript + Vite, Tailwind CSS + shadcn/ui, React Router v6,
%         TanStack React Query
%   \item \textbf{Backend:} FastAPI (Python), structured as a layered monolith
%         (\texttt{routes $\to$ schemas $\to$ services $\to$ repositories}); a microservices
%         decomposition is prototyped on the \texttt{comparison} branch
%   \item \textbf{Database:} Firebase Firestore (NoSQL, per-user partitioned)
%   \item \textbf{Authentication:} Firebase Authentication (email/password + Google OAuth)
%   \item \textbf{Image analysis:} ONNX Runtime with a fine-tuned MobileNet-family model,
%         plus a tiered nutrition lookup chain
%   \item \textbf{Observability:} In-process Prometheus-compatible metrics, structured JSON logging,
%         correlation-ID middleware
%   \item \textbf{Testing:} pytest (backend, 56 tests, 76\% line coverage) and Vitest (frontend)
%   \item \textbf{Deployment:} Docker + Docker Compose (gateway + 6 services, target architecture)
% \end{itemize}

% \subsection{Implemented features (current repository state)}
% \begin{itemize}[leftmargin=1.2em]
%   \item Authentication (email/password + Google OAuth), protected routes
%   \item Meal logging (manual) with multi-item nutrition tracking
%   \item Image-based meal analysis to auto-fill detected foods and estimated nutrition
%   \item Workout logging (type, duration, intensity, calories burned, notes)
%   \item Daily goals (calories/macros, optional target weight)
%   \item Unified daily dashboard summary (calories in/out/net, macro breakdown, counts)
%   \item History view for meals/workouts with delete
%   \item User profile (personal details + dietary preferences/allergies)
%   \item Progress analytics for a date range (trend charts + breakdowns)
%   \item Workout plans (admin-managed plans, publicly accessible)
%   \item Personalized workout recommendations (rule-based engine, goal- and history-aware)
% \end{itemize}

% ===========================
% Task 1: Requirements and Subsystems
% ===========================
\section{Task 1: Requirements and Subsystems}

\subsection{Functional requirements}

The initial proposal defined the core functional scope: account management, meal logging, workout tracking, goals, dashboard, analytics, and recommendations.
During implementation, several flows were expanded and clarified to make the system usable end-to-end.

\begin{longtable}{p{1.2cm} p{5.8cm} p{6.5cm}}
  \toprule
  \textbf{ID} & \textbf{Requirement} & \textbf{Success criteria} \\
  \midrule
  \endhead

  FR-1  & Users shall be able to register, login/logout, and maintain an authenticated session
        & \texttt{POST /auth/register} returns 201; \texttt{POST /auth/login} returns a valid
          Firebase ID token; \texttt{POST /auth/logout} clears the client token; any
          \texttt{/api/} call without a valid token returns HTTP 401. \\[2pt]

  FR-2  & Users shall be able to sign in using email/password and Google OAuth
        & Both providers issue a Firebase ID token accepted by \texttt{get\_current\_user};
          neither path exposes any credential in the frontend bundle or URL. \\[2pt]

  FR-3  & Authenticated users shall be able to create and update a personal profile
        & \texttt{GET /api/v1/profile} returns all stored fields; \texttt{PUT /api/v1/profile}
          with changed fields returns 200 and the updated document; subsequent GET reflects
          the change. \\[2pt]

  FR-4  & Authenticated users shall be able to log a meal with meal type, date, time, and
          one or more food items
        & \texttt{POST /api/v1/meals} with a valid payload returns 201; the meal is
          retrievable via \texttt{GET /api/v1/meals?dateFrom=\ldots}; a request with
          zero items or a missing required field returns 422. \\[2pt]

  FR-5  & Each meal item shall include name, quantity, unit, calories, protein, carbs, and fat
        & All seven fields present and non-null; quantity $> 0$; calories, protein, carbs,
          fat $\geq 0$; Pydantic validation rejects out-of-range values with 422. \\[2pt]

  FR-6  & The system shall calculate and persist meal totals server-side
        & Totals in the stored meal document equal the server-computed sum across items;
          submitting a payload with manually inflated client totals has no effect on
          the persisted values. \\[2pt]

  FR-7  & Users shall be able to upload a meal photo for automatic food detection
          and nutrition autofill
        & \texttt{POST /api/v1/meals/analyze-image} with a valid image returns at least one
          detected food item with non-zero macro estimates; an unrecognisable image returns
          a structured error (\texttt{no\_food\_detected}), not a 500. \\[2pt]

  FR-8  & Authenticated users shall be able to log workouts
        & \texttt{POST /api/v1/workouts} returns 201; type, duration ($> 0$), intensity,
          calories burned ($\geq 0$) are stored; missing required fields return 422. \\[2pt]

  FR-9  & Users shall be able to set and update daily goals
        & \texttt{PUT /api/v1/goals} creates or replaces the active goal; only one active
          goal exists per user at any time; \texttt{GET /api/v1/goals} returns it. \\[2pt]

  FR-10 & Users shall be able to view a unified daily dashboard
        & \texttt{GET /api/v1/dashboard?date=\ldots} returns \texttt{calories\_in},
          \texttt{calories\_out}, \texttt{net\_calories}, macro breakdown, meal count, and
          workout count; response includes all fields even when no data exists for that day
          (zeroed values, not missing keys). \\[2pt]

  FR-11 & Users shall be able to view history of logged meals and workouts and delete entries
        & History endpoint returns items for the requested date range; \texttt{DELETE
          /api/v1/meals/\{id\}} returns 204 and the entry is absent from subsequent history
          calls; deleting a non-existent ID returns 404. \\[2pt]

  FR-12 & Users shall be able to view progress analytics across a date range
        & \texttt{GET /api/v1/analytics?dateFrom=\ldots\&dateTo=\ldots} returns trend
          data for every day in the range; requesting a future range returns an empty
          dataset, not an error. \\[2pt]

  FR-13 & The system shall provide personalized workout recommendations based on goals
          and historical behavior
        & \texttt{GET /api/v1/recommendations} returns $\geq 1$ recommendation when the
          user has an active goal; the recommendation strategy matches the goal type
          (strength / cardio / weight-loss). \\

  \bottomrule
\end{longtable}


% \subsection{Persistence model (Firestore)}
% All data is stored per user under a \texttt{users/\{uid\}} partition. The core Firestore paths:
% \begin{lstlisting}
% users/{uid}/meals/{mealId}          <- meal summary + denormalized items[]
% users/{uid}/workouts/{workoutId}
% users/{uid}/goals/active
% users/{uid}/profile/main
% users/{uid}/recommendations/{id}
% plans/{planId}                      <- admin-managed plans (global collection)
% \end{lstlisting}

% \textbf{N+1 optimization:} meal items are denormalized into the parent meal document at write time,
% eliminating the sub-collection fan-out on reads. Listing $N$ meals now costs $1$ Firestore
% collection read rather than the former $N{+}1$ (see Section~\ref{sec:nfr-quantification}).

\subsection{Non-functional requirements}

The proposal identified six quality attributes: scalability, performance, extensibility,
maintainability, security/privacy, and reliability/testability.
We split Reliability and Testability into separate rows and add Data Integrity and
Observability based on implementation evidence.

\begin{longtable}{p{1.2cm} p{2.5cm} p{4.8cm} p{5.0cm}}
  \toprule
  \textbf{ID} & \textbf{Attribute} & \textbf{Necessity} & \textbf{Success threshold} \\
  \midrule
  \endhead

  NFR-1 & Security / Privacy
        & All \texttt{/api/} endpoints reject unauthenticated requests;
          reads/writes scoped to the authenticated UID only.
        & 0 endpoints accessible without a valid Firebase ID token;
          \texttt{gitleaks}: 0 secrets in HEAD;
          \texttt{slowapi} rate limits active on auth and image-analysis routes. \\[2pt]

  NFR-2 & Data Integrity
        & Business-logic aggregates computed server-side; input validated at API boundary.
        & 0 client-trusted aggregate values in any \texttt{POST} schema;
          $\geq$87\% of input fields carry explicit Pydantic range / format constraints. \\[2pt]

  NFR-3 & Performance
        & Firestore reads sub-linear in history size;
          initial JS bundle small for fast initial load.
        & $\leq 1$ Firestore read per meal-list page (N+1 fix);
          initial JS bundle $\leq 150$\,KB gzip
          (measured: \textbf{137.14\,KB}). \\[2pt]

  NFR-4 & Scalability
        & Multiple users and increasing per-user data supported without correctness
          issues; read-heavy endpoints cached.
        & No per-user data leakage under concurrent load;
          routing layer sustains $\geq$70\,RPS with p95 $\leq$10\,ms
          (measured: \textbf{73.8\,RPS, p95 3.78\,ms}). \\[2pt]

  NFR-5 & Maintainability
        & Strict separation of concerns; cyclomatic complexity bounded.
        & Average CC $\leq$ grade B (radon);
          no single function $>$ grade D;
          measured: \textbf{avg A (2.98)}, one outlier at D (27). \\[2pt]

  NFR-6 & Extensibility
        & Adding a new health data source requires no edits to existing modules.
        & Extension cost $=$ \textbf{+1 class, 0 edits};
          \texttt{test\_open\_closed\_extension\_via\_new\_provider}
          passes in CI. \\[2pt]

  NFR-7 & Reliability
        & Graceful degradation under partial failures;
          health probes exposed.
        & Partial-failure response (non-5xx) on dashboard when one domain is
          degraded; \texttt{/health} and \texttt{/ready} endpoints return 200 / 503
          correctly. \\[2pt]

  NFR-8 & Testability
        & Core logic unit-testable; CI enforces a coverage floor.
        & Backend line coverage $\geq$60\% (CI gate \texttt{--cov-fail-under=60});
          measured: \textbf{76\%}, \textbf{56/56} tests passing in \textbf{\textasciitilde{}1\,s}. \\[2pt]

  NFR-9 & Observability
        & Every request carries a correlation ID;
          Prometheus metrics exposed.
        & 100\% of requests carry \texttt{x-request-id} by construction;
          \texttt{/metrics} exposes $\geq$2 metric families
          (measured: \textbf{2}: counter + histogram). \\

  \bottomrule
\end{longtable}

\subsection{Key architecturally significant requirements (ASRs)}
A requirement is architecturally significant when satisfying it forces a structural decision that
cannot be easily refactored later, rules out otherwise attractive alternatives, or has
system-wide ripple effects. The six below meet that bar.

\begin{enumerate}[leftmargin=1.4em]

  \item \textbf{Authentication + per-user data isolation} (FR-1/FR-2, NFR-1) \\
  \textit{Why it is a driver:} Health data is private; a data-isolation failure is a security
  incident, not a bug. Allowing direct client-to-Firestore writes would scatter auth logic
  across the frontend and require maintaining Firestore security rules in sync with business
  rules. \\
  \textit{Structural decision forced:} All writes go through a backend that verifies the
  Firebase ID token and uses the extracted \texttt{uid} to scope every Firestore path
  (\texttt{users/\{uid\}/...}). The \texttt{get\_current\_user} FastAPI dependency is wired
  into every route via \texttt{Depends}, making accidental unauthenticated access a compile-time
  mistake rather than a runtime gap.

  \item \textbf{Server-side totals and Pydantic validation} (FR-4/5/6, NFR-2) \\
  \textit{Why it is a driver:} If the backend accepted client-computed nutrition totals, a
  tampered client could store arbitrary macro values. Pydantic validation at the route boundary
  is the only place guaranteed to run regardless of client behaviour. \\
  \textit{Structural decision forced:} Business logic (total computation) must live in the
  service layer, not in routes or repositories. This drove the four-layer architecture; removing
  the service layer would require pushing logic into routes (breaking testability) or trusting
  clients (breaking integrity).

  \item \textbf{ML inference on the hot request path} (FR-7, NFR-3) \\
  \textit{Why it is a driver:} ONNX inference is CPU-bound and takes hundreds of milliseconds.
  Running it synchronously on FastAPI's single-threaded async event loop would block every
  concurrent request---all other users would see elevated latency during any image upload. \\
  \textit{Structural decision forced:} Inference is offloaded via \texttt{asyncio.to\_thread},
  introducing an explicit concurrency seam in the meals service. The nutrition lookup chain
  (Strategy + Chain-of-Responsibility) was extracted from the service because adding new food
  databases would otherwise require editing inference-adjacent code.

  \item \textbf{Extensible nutrition lookup} (NFR-6) \\
  \textit{Why it is a driver:} Nutrition accuracy is a product quality concern that will improve
  over time (adding regional databases, external APIs). Hard-coding the lookup order inside a
  service file means each improvement is a risky diff on business-logic code. \\
  \textit{Structural decision forced:} The \texttt{NutritionProvider} protocol and
  \texttt{NutritionResolver} chain (ADR-0004). Any design that keeps the lookup inside a single
  function fails NFR-6 because the extension cost is always $\geq 1$ edit.

  \item \textbf{Cache-aside for read-heavy aggregation endpoints} (FR-10, NFR-4) \\
  \textit{Why it is a driver:} The dashboard and analytics services re-read every meal and
  workout for the user on every request. Under any real load this becomes a Firestore read
  amplification problem. Without a cache, horizontal scaling only moves the bottleneck to
  Firestore quotas. \\
  \textit{Structural decision forced:} An explicit \texttt{Cache} protocol with pluggable backends
  (in-memory for the monolith, Redis-ready for multi-instance). Services depend on the protocol,
  not a concrete backend---Dependency Inversion Principle applied at the infrastructure boundary.

  \item \textbf{Progress analytics read model} (FR-12, NFR-3/4) \\
  \textit{Why it is a driver:} Analytics queries span arbitrary date ranges over unbounded
  per-user history. A naive scan-and-aggregate on every request grows linearly with user
  data volume, violating NFR-3. \\
  \textit{Structural decision forced:} Firestore predicate push-down (\texttt{where} clauses)
  and write-time denormalization (items embedded in meal documents, ADR-0003) so that list
  queries never fan out into per-document sub-collection reads. Together these keep the
  analytics service read cost at $O(\mathrm{range\ size})$ rather than $O(\mathrm{total\ history})$.

\end{enumerate}

\subsection{Subsystem overview}

\textbf{Decomposition principle.}
Subsystems are grouped by \emph{data ownership and capability boundary}, not by technical layer.
Each domain service owns one Firestore collection namespace and exposes one cohesive set of
operations on it; no two services read or write the same collection.
The frontend is separated because it runs in an untrusted execution environment (the browser)
and must never hold service-account credentials.
Image recognition and nutrition lookup are extracted because they carry their own heavyweight
dependencies (ONNX runtime, food databases) whose failure modes and extension cost differ
from plain CRUD, and they are only consumed by the meals pipeline.

\textbf{Cross-cutting interfaces.}
Three concerns span every subsystem and are implemented once rather than replicated:
\begin{itemize}[leftmargin=1.2em, topsep=2pt]
  \item \textbf{Authentication} --- \texttt{get\_current\_user} (auth subsystem) is injected
        via FastAPI \texttt{Depends} into every domain route; no domain service re-implements
        token verification.
  \item \textbf{Error propagation} --- all domain errors are subclasses of \texttt{DomainError}
        (errors subsystem); a single handler in \texttt{main.py} maps them to HTTP status codes,
        so domain services never import from the transport layer.
  \item \textbf{Observability} --- correlation-ID middleware, structured logging, and Prometheus
        metrics are registered once on the FastAPI app; every request is automatically traced
        and timed regardless of which domain route handles it.
\end{itemize}

\begin{longtable}{p{3.5cm} p{5.5cm} p{5.2cm}}
  \toprule
  \textbf{Subsystem} & \textbf{Role} & \textbf{Interfaces with} \\
  \midrule
  \endhead

  Web frontend
  & All user-facing interactions: auth flows, dashboard, logging, analytics, recommendations, plans.
  & Authentication (token acquisition); all backend API routes (via \texttt{apiFetch}). \\

  Authentication
  & Firebase ID token issuance (client SDK) and verification (Admin SDK via
    \texttt{get\_current\_user}).
  & All domain services (cross-cutting); Firebase Auth cloud service. \\

  Meals service
  & CRUD for meal records; delegates image analysis and nutrition lookup.
  & Image recognition; nutrition lookup; persistence; authentication. \\

  Workouts service
  & CRUD for workout records.
  & Persistence; authentication. \\

  Goals service
  & Manage the single active goal per user.
  & Persistence; authentication; recommendations (reads active goal). \\

  Dashboard service
  & Aggregate daily calories-in/out, macro breakdown, counts; cached.
  & Persistence; caching subsystem; authentication. \\

  Profile service
  & Store and retrieve user profile (demographics, dietary preferences).
  & Persistence; authentication; recommendations (reads profile). \\

  Analytics service
  & Compute trend metrics over a date range; predicate push-down to Firestore.
  & Persistence; caching subsystem; authentication. \\

  Recommendations
  & Rule-based workout plan generation using goal + history; Strategy pattern.
  & Goals service; profile service; persistence; authentication. \\

  Plans service
  & Admin-managed global workout plan catalogue; admin auth via custom header.
  & Persistence; authentication. \\

  Image recognition
  & ONNX MobileNet food classification; async thread-pool offload.
  & Meals service (internal call only); ONNX runtime (file dependency). \\

  Nutrition lookup
  & Tiered food-name-to-macro resolution via Chain-of-Responsibility.
  & Meals service (internal call only); food database files. \\

  Persistence
  & Firebase Firestore read/write under \texttt{users/\{uid\}/} and \texttt{plans/}.
  & All domain services (cross-cutting); Firebase Firestore cloud service. \\

  Observability
  & Prometheus metrics, structured JSON logging, correlation-ID propagation.
  & All subsystems (cross-cutting, registered once on the FastAPI app). \\

  \bottomrule
\end{longtable}

% ===========================
% Task 2: Architecture Framework
% ===========================
\newpage
\section{Task 2: Architecture Framework}

\subsection{Stakeholders}

\begin{longtable}{p{5.0cm} p{8.5cm}}
  \toprule
  \textbf{Stakeholder} & \textbf{Primary concerns} \\
  \midrule
  \endhead

  End user (health tracker)
    & Fast entry, accurate nutrition totals, trend visibility,
      private health data, user-friendly UI. \\
  Plan Admin
    & Clear content management paths, reliable data integrity,
      and a straightforward way to add/update workout plans. \\
  Development team
    & Clear module boundaries, local dev workflow, testable
      modules, straightforward extension paths. \\
  \bottomrule
\end{longtable}

\subsection{Architecture Views}

\subsubsection{Architectural Overview}

\begin{figure}[H]
  \centering
  \includegraphics[width=0.70\linewidth]{"../docs/Diagrams/layered"}
  \caption{High-level architectural overview of \ProjectName.}
\end{figure}

\subsubsection{API Design View}

\begin{figure}[H]
  \centering
  \includegraphics[width=0.88\linewidth]{"../docs/Diagrams/API Design View"}
  \caption{REST API design view}
\end{figure}

\subsubsection{Frontend and Backend Structure}

\begin{figure}[H]
  \centering
  \includegraphics[width=0.50\linewidth]{"../docs/Diagrams/Frontend_Backend"}
  \caption{Component/layer diagram}
\end{figure}

\subsubsection{Frontend C3 Diagram}

\begin{figure}[H]
  \centering
  \includegraphics[width=\linewidth]{"../docs/Diagrams/frontend_c3"}
  \caption{C3 diagram for the frontend}
\end{figure}

\subsubsection{Backend C3 Diagram}

\begin{figure}[H]
  \centering
  \includegraphics[width=\linewidth]{"../docs/Diagrams/backend_c3"}
  \caption{C3 diagram for the backend}
\end{figure}

\subsubsection{Interaction Flow}

\begin{figure}[H]
  \centering
  \includegraphics[width=0.82\linewidth]{"../docs/Diagrams/Interaction Flow"}
  \caption{End-to-end interaction flow}
\end{figure}

\subsubsection{Meal Logging Sequence}

\begin{figure}[H]
  \centering
  \includegraphics[width=\linewidth]{"../docs/Diagrams/Meal Logging"}
  \caption{UML sequence diagram for the Log Meal flow}
\end{figure}

\subsubsection{Workout Logging Sequence}

\begin{figure}[H]
  \centering
  \includegraphics[width=\linewidth]{"../docs/Diagrams/Workout Sequence"}
  \caption{UML sequence diagram for the Log Workout flow}
\end{figure}

% ===========================
% Major Design Decisions (ADRs)
% ===========================
\subsection{Major Design Decisions (Architecture Decision Records)}


\subsubsection{ADR-0001: Strict layered monolith architecture}
\textbf{Status:} Accepted (layered monolith) \\
\textbf{Context:} \ProjectName\ spans multiple feature domains (meals, workouts, goals, dashboard,
profile, analytics, recommendations, plans). Early development called for clear module ownership and easy local setup, while the proposal simultaneously required
demonstrating a decomposable service boundary. Premature distribution would have added
inter-service networking and distributed tracing overhead before the domain models had stabilised. \\
\textbf{Decision:} Deliver the prototype as a single FastAPI application organised in a strict
\textit{routes $\to$ schemas $\to$ services $\to$ repositories} layering. Each layer may only import from the layer directly below it. \\
\textbf{Consequences:}
\begin{itemize}[leftmargin=1.2em]
  \item (+) Single process: one port, one log stream, one-command start.
  \item (+) Layer discipline means each domain can be extracted to its own service with changes
        confined to one layer.
  \item (--) A process-level bug (e.g., OOM, unhandled exception) can affect all feature domains.
        Mitigated by the domain error hierarchy (see ADR-0002) and layer isolation.
  \item (--) Horizontal scaling requires scaling the entire process; individual aspects
        (e.g., image analysis) cannot be scaled independently in the monolith.
\end{itemize}

\subsubsection{ADR-0002: Authenticated backend API boundary for all user data writes}
\textbf{Status:} Accepted \\
\textbf{Context:} Meal, workout, goal, and profile records are private health data. The alternative
of direct client-side Firestore writes would require maintaining Firestore security rules in sync
with client logic and would allow a tampered client to forge nutrition totals or access another
user's partition. \\
\textbf{Decision:} All data mutations pass through FastAPI route handlers. The frontend obtains a
short-lived Firebase ID token and \texttt{apiFetch} attaches it as a Bearer header. The backend
\texttt{get\_current\_user} dependency (in the auth/security subsystem) verifies the token via the
Firebase Admin SDK and extracts the verified \texttt{uid}, which is the only identity used to
scope Firestore reads and writes. The client SDK is used for authentication only; it never touches
Firestore directly.
diagrams. \\
\textbf{Consequences:}
\begin{itemize}[leftmargin=1.2em]
  \item (+) Auth, validation, and logging are enforced at one trusted point; routes cannot
        accidentally bypass them (FastAPI \texttt{Depends} wiring makes omission a type error).
  \item (+) Server-side nutrition total computation eliminates the risk of a manipulated client
        payload - a key requirement for data integrity (NFR-2).
  \item (+) Service layer is testable without a live Firestore instance (uid can be a plain string
        in tests).
  \item (--) Prototype requires both frontend and backend servers running concurrently.
  \item (--) One additional network hop compared to direct Firestore client writes.
\end{itemize}

\subsubsection{ADR-0003: Denormalised Firestore meal document}
\textbf{Status:} Accepted \\
\textbf{Context:} The initial Firestore model stored meal items in a sub-collection. Listing a user's meal history
therefore required $N{+}1$ reads: one collection scan to get $N$ meal summaries, plus one
sub-collection read per meal to get items. At page size 100, this was 101 Firestore reads per
list call - a direct violation of NFR-3 (sub-linear read cost). \\
\textbf{Decision:} Embed food items as an \texttt{items[]} array directly in the parent meal
document at write time, alongside server-computed aggregate totals. Meal summary and item detail
are now co-located in a single document. A backward-compatibility read path in the meals
service handles any legacy sub-collection rows. \\
\textbf{Consequences:}
\begin{itemize}[leftmargin=1.2em]
  \item (+) Listing $N$ meals costs exactly \textbf{1 Firestore read} instead of $N{+}1$, a $100\times$ reduction at $N{=}100$.
  \item (+) Removes the async fan-out; the service layer simplifies from a gather step to a single collection scan.
  \item (--) Meal documents grow with item count; the 1 MB Firestore document limit bounds the maximum items per meal.
\end{itemize}

\subsubsection{ADR-0004: Strategy + Chain-of-Responsibility for nutrition lookup}
\textbf{Status:} Accepted \\
\textbf{Context:} The image analysis endpoint must map a detected food name (e.g., ``samosa'',
``banana'') to calorie and macro estimates. An initial design hard-coded the priority order and
fallback logic inside the meals service. Adding a new food database required editing that file,
violating the Open/Closed Principle and making nutrition accuracy improvements a risky diff
on a business-logic file. \\
\textbf{Decision:} Extract the lookup into a \texttt{NutritionProvider} protocol with a single
\texttt{lookup(name)} method, and a \texttt{NutritionResolver} that walks a prioritised chain
of providers, short-circuiting on the first non-\texttt{None} hit. Three providers ship:
an Indian-cuisine database (highest priority), a global common-foods database, and a
\texttt{DefaultEstimateProvider} as a terminal fallback. The chain is wired once at startup
in the nutrition subsystem; new sources require only a new class and an insertion point---no
edits to existing providers or the resolver. \\
\textbf{Consequences:}
\begin{itemize}[leftmargin=1.2em]
  \item (+) Extension cost is just adding a new class.
  \item (+) Chain short-circuits on the first hit; the expensive \texttt{DefaultEstimateProvider} is only reached for unknown foods.
  \item (--) Chain order is configuration, not type-system enforced; a misconfigured priority order silently returns lower-accuracy data.
\end{itemize}

\newpage
\subsubsection{ADR-0005: Graceful dashboard degradation under partial service failures}
\textbf{Status:} Accepted \\
\textbf{Context:} The dashboard aggregates meals, workouts, goals, and profile data in a single
response. In a multi-service architecture, if any upstream times out, a naive synchronous
gather will block the entire dashboard until the slowest dependency responds or times out.
Even in the monolith, a slow Firestore read for one domain should not stall the response for
the others. Users should see a partial but correct summary rather than a full failure
(NFR-7). \\
\textbf{Decision:} In the layered monolith: the dashboard service queries each domain independently
with short timeouts; missing data produces a zero-valued field rather than an exception. The
cache-aside layer (in the caching subsystem) absorbs repeated reads: dashboard responses are
cached with a TTL so a slow upstream affects only the first uncached request. \\
\textbf{Consequences:}
\begin{itemize}[leftmargin=1.2em]
  \item (+) Users can still view available health data even when one domain's backend is degraded.
  \item (+) Cache-aside (TTL-based) prevents one slow Firestore read from blocking every user
        request on the hot dashboard endpoint.
  \item (--) Partial responses require the frontend to handle missing fields gracefully.
\end{itemize}

% ===========================
% Task 3: Architectural Tactics and Patterns
% ===========================
\newpage
\section{Task 3: Architectural Tactics and Patterns}

\subsection{Architectural tactics}

\begin{longtable}{p{3.5cm} p{2.4cm} p{4.2cm} p{4.0cm}}
  \toprule
  \textbf{Tactic} & \textbf{Quality attr.} & \textbf{How it is applied} & \textbf{Where in code} \\
  \midrule
  \endhead

  Authenticate at API boundary
    & Security
    & \texttt{get\_current\_user} verifies the Firebase ID token; any missing or invalid
      token raises \texttt{AuthorizationError} which maps to HTTP 401 before the route
      handler body executes.
    & \texttt{app/config/security.py}; injected via \texttt{Depends(get\_current\_user)}
      in every route file under \texttt{app/routes/}. \\[2pt]

  Validate at multiple boundaries
    & Integrity
    & Frontend client-side validation blocks submission for UX.
      Backend Pydantic schemas enforce non-negative values, max lengths, regex patterns,
      and bounded list sizes as the authoritative constraint.
    & \texttt{app/schemas/*.py} (backend); \texttt{frontend/src/pages/LogMeal.tsx} and
      sibling pages (frontend form validation). \\[2pt]

  Server-side aggregate computation
    & Integrity
    & Nutrition totals and dashboard aggregates are computed in the service layer on
      every write; no endpoint accepts pre-computed totals from the client.
    & \texttt{app/services/meals\_service.py} (nutrition totals at write time);
      \texttt{app/services/dashboard\_service.py} (daily aggregates). \\[2pt]

  Cache-Aside with pluggable backends
    & Performance / Scalability
    & Dashboard and analytics responses are stored under a keyed TTL cache;
      a cache miss falls through to Firestore; invalidation happens at write time.
    & \texttt{app/config/cache.py} (\texttt{Cache} protocol, \texttt{InMemoryTTLCache},
      \texttt{RedisCache}); used in dashboard and analytics services via
      \texttt{Depends(get\_cache)}. \\[2pt]

  Async thread-pool offload for CPU-bound I/O
    & Performance
    & ONNX image inference is dispatched to a thread-pool worker so the FastAPI
      async event loop is never blocked; other requests are served concurrently
      during inference.
    & \texttt{app/services/meals\_service.py}:
      \texttt{await asyncio.to\_thread(recognizer.predict, image\_bytes)}. \\[2pt]

  Correlation-ID + structured logging
    & Observability
    & HTTP middleware mints a UUID4 \texttt{x-request-id} (or accepts an upstream one),
      attaches it to every log line, and echoes it on the response header. In microservices
      mode the gateway propagates the same ID to upstreams.
    & \texttt{app/main.py} (middleware, lines 55--90);
      \texttt{microservices/\_base.py} + \texttt{microservices/gateway/proxy.py}
      (\texttt{comparison} branch). \\[2pt]

  Prometheus in-process metrics
    & Observability
    & A zero-dependency metrics registry counts requests and records latency
      histograms (11 buckets, 5\,ms--10\,s) per route; \texttt{/metrics} exposes
      Prometheus text format without requiring an external agent.
    & \texttt{app/observability/metrics.py};
      registered via \texttt{register\_metrics\_middleware(app)} and
      \texttt{register\_metrics\_endpoint(app)} at the bottom of \texttt{app/main.py}. \\[2pt]

  React Query stale-while-revalidate + mutation invalidation
    & Consistency
    & After every successful mutation (meal log, workout log, goal update) the frontend
      invalidates the affected query keys, triggering a background refetch so the
      dashboard and history views reflect the new server state immediately.
    & \texttt{frontend/src/hooks/useApi.ts};
      mutation callbacks call \texttt{queryClient.invalidateQueries(\{queryKey: [\ldots]\})}
      targeting dashboard and history keys. \\

  \bottomrule
\end{longtable}

\subsection{Implementation patterns}

\subsubsection{Pattern 1: Strategy Pattern (Recommendation Engine)}

\textbf{Intent:} Define a family of algorithms, encapsulate each one, and make them
interchangeable. The pattern lets the algorithm vary independently from clients that use it.

\textbf{Role in \ProjectName:} The \texttt{RecommendationEngine} selects and executes a
workout-plan-building algorithm based on the user's goals and profile, without the
engine itself knowing which concrete strategy is active.

\textbf{Structure:}
\begin{itemize}[leftmargin=1.2em]
  \item \textbf{Context} --- \texttt{RecommendationEngine}
        (\texttt{backend/app/services/recommendation\_engine.py}): holds a reference to a
        \texttt{RecommendationStrategy} and delegates generation to it. The private \texttt{resolve\_strategy(context)}
        method selects the concrete strategy based on the user's stated goal.
  \item \textbf{Strategy interface} --- \texttt{RecommendationStrategy}: declares build\_plan(context) and build\_week\_plan(context).
  \item \textbf{Concrete strategies:}
    \begin{itemize}[leftmargin=1.2em]
      \item \texttt{StrengthTrainingStrategy} --- generates strength-focused plans
      \item \texttt{CardioEnduranceStrategy} --- generates cardio/endurance plans
      \item \texttt{WeightLossStrategy} --- generates calorie-deficit-oriented plans
    \end{itemize}
\end{itemize}

\begin{figure}[H]
  \centering
  \includegraphics[width=0.72\linewidth]{"../docs/Diagrams/Strategy Pattern_Recommendation"}
  \caption{Strategy Pattern applied to the Recommendation Engine.}
\end{figure}

\textbf{How it addresses NFRs:}
\begin{itemize}[leftmargin=1.2em]
  \item \textbf{Extensibility:} adding a new goal type (e.g., \textit{Flexibility}) requires only a new class implementing RecommendationStrategy.
  \item \textbf{Maintainability:} each strategy is a self-contained module with a single responsibility; they can be read, tested, and evolved independently.
  \item \textbf{Testability:} the strategy interface can be mocked in unit tests RecommendationEngine without instantiating any real plan-building logic.
\end{itemize}


% ---

\subsubsection{Pattern 2: Repository Pattern}

\textbf{Role in \ProjectName:} All Firestore interactions are channelled through repository classes. Service code never constructs Firestore document paths or executes queries directly, making the data access layer independently testable and swappable.

\textbf{Structure:}
\begin{itemize}[leftmargin=1.2em]
  \item \textbf{Abstract base} --- \texttt{BaseRepository}: declares the common interface
        \texttt{get\_by\_id(id)}, \texttt{add(document)}, and \texttt{update(id, document)}.
  \item \textbf{Concrete repositories:}
    \begin{itemize}[leftmargin=1.2em]
      \item \texttt{MealsRepository} --- adds \texttt{get\_meals\_by\_date(user\_id, date)};
            encapsulates the Firestore path and the date-filter predicate.
      \item \texttt{WorkoutsRepository} --- adds
            \texttt{get\_workouts\_in\_range(user\_id, start\_date, end\_date)}; encapsulates the Firestore path and range query.
    \end{itemize}
  \item Analogous repositories exist for Goals, Profile, Plans, and Recommendations.
\end{itemize}

\begin{figure}[H]
  \centering
  \includegraphics[width=0.78\linewidth]{"../docs/Diagrams/Repository Pattern"}
  \caption{Repository Pattern: \texttt{BaseRepository} defines the common data-access interface;
    \texttt{MealsRepository} and \texttt{WorkoutsRepository} extend it with domain-specific
    query methods, hiding all Firestore path logic from the service layer.}
\end{figure}

\textbf{How it addresses NFRs:}
\begin{itemize}[leftmargin=1.2em]
  \item \textbf{Testability:} Easy to mock the actual db implementation; service tests never touch real Firestore.
  \item \textbf{Maintainability:} Firestore path construction is in one place; renaming a collection requires editing only the relevant repository file.
  \item \textbf{Extensibility:} adding a new collection (e.g., weight logs) means creating a new repository subclass without touching any existing service.
\end{itemize}

% ---

\subsubsection*{Other design patterns present}

\begin{itemize}[leftmargin=1.2em]

  \item \textbf{Chain of Responsibility}
        (\texttt{backend/app/services/nutrition/providers.py}) \\
        \texttt{NutritionResolver} walks a list of \texttt{NutritionProvider} objects in
        priority order and short-circuits on the first non-\texttt{None} hit
        (Indian foods $\to$ global foods $\to$ default estimate). Each provider decides
        independently whether it can handle the request; none has knowledge of its neighbours.

  \item \textbf{Facade}
        (\texttt{backend/app/services/*.py}) \\
        Each service module (e.g., \texttt{meals\_service.py}, \texttt{dashboard\_service.py})
        acts as a Facade: it presents a simple, high-level interface to the route layer
        (\texttt{log\_meal()}, \texttt{get\_dashboard()}) while internally coordinating the
        repository, cache, image recogniser, and domain error hierarchy. The route handler
        never needs to know about Firestore or ONNX.

  \item \textbf{Observer}
        (\texttt{frontend/src/} --- TanStack React Query) \\
        After every successful mutation (meal log, workout log, goal update), the frontend ensures React Query re-fetches all registered observers of the invalidated key, ensuring the dashboard and history views automatically reflect the latest server state without manual re-rendering.

\end{itemize}

% ===========================
% Task 4: Prototype Implementation and Architecture Analysis
% ===========================
\newpage
\section{Task 4: Prototype Implementation and Architecture Analysis}

\subsection{Implemented prototype}
\subsubsection*{Today's Overview on the Dashboard}
\begin{figure}[H]
  \centering
  \includegraphics[width=0.72\linewidth]{"../docs/UI/Dashboard.png"}
  \caption{Dashboard UI with Today's Overview.}
\end{figure}

\subsubsection*{Log Meal form with image analysis}
\begin{figure}[H]
  \centering
  \includegraphics[width=0.72\linewidth]{"../docs/UI/Log-Meal.png"}
  \caption{Log Meal form with image analysis results.}
\end{figure}  

\subsubsection*{Profile page with editable fields}
\begin{figure}[H]
  \centering
  \includegraphics[width=0.72\linewidth]{"../docs/UI/Profile.png"}
  \caption{Profile page with editable fields.}
\end{figure}

\subsubsection*{History page}
\begin{figure}[H]
  \centering
  \includegraphics[width=0.72\linewidth]{"../docs/UI/History.png"}
  \caption{History page showing past workouts.}
\end{figure}

\subsubsection*{Workout recommendation page}
\begin{figure}[H]
  \centering
  \includegraphics[width=0.72\linewidth]{"../docs/UI/Plans.png"}
  \caption{Workout recommendation page showing available plans.}
\end{figure} 

\subsubsection*{Log Workout form}
\begin{figure}[H]
  \centering  
  \includegraphics[width=0.72\linewidth]{"../docs/UI/Log-Workout.png"}
  \caption{Log Workout form.}
\end{figure}

\subsubsection*{Progress view with charts}
\begin{figure}[H]
  \centering
  \includegraphics[width=0.72\linewidth]{"../docs/UI/Progress.png"}
  \caption{Progress view with charts.}
\end{figure}


\subsection{Architecture comparison: layered monolith vs.\ microservices}

The implemented architecture (\texttt{layered} branch) is a \textbf{strict layered monolith}.
The comparison architecture (\texttt{comparison} branch) decomposes the same feature set into
\textbf{six independently-runnable FastAPI services} behind a gateway.

\begin{longtable}{p{2.8cm} p{6.3cm} p{6.1cm}}
  \toprule
  \textbf{Concern} & \textbf{Layered monolith (implemented)} & \textbf{Microservices (comparison branch)} \\
  \midrule
  \endhead

  Security
    & Single token-verification dependency; one CORS configuration; one TLS endpoint.
    & Each service independently verifies tokens (same Firebase Admin SDK); gateway owns CORS; larger attack surface but more granular access control. \\

  Data integrity
    & Single process; service-to-service calls are in-process function calls with no serialisation overhead.
    & Service-to-service calls cross network boundaries; JSON serialisation adds latency but makes contracts explicit. \\

  Maintainability
    & Strict layer separation is enforced by module imports; changes to one feature area can still accidentally reach another via shared utilities.
    & Physical process isolation prevents cross-domain import accidents; each service has its own \texttt{requirements.txt} and deployment lifecycle. \\

  Performance
    & No inter-service network hops. Measured: mean \textbf{2.84 ms}, p95 \textbf{4.08 ms}
      on \texttt{GET /api/v1/meals} (80 RPS, 15 s bench). Cache-aside is per-process (in-memory).
    & One gateway $\to$ upstream HTTP hop adds \textbf{+1.27 ms mean (+45\%)}, p95 \textbf{5.32 ms}
      on the same route. Cache must be shared (Redis) across service instances. \\

  Scalability
    & \textbf{73.8 RPS} (p95 3.78 ms) at the routing layer (\texttt{/health}, 0 errors).
      Scaling requires scaling the entire process; a single hot path (e.g., image analysis)
      cannot scale independently.
    & \textbf{73.9 RPS} (p95 2.73 ms) at the gateway on the same route.
      Individual services can be scaled horizontally per load profile
      (e.g., scale the analytics or image-analysis service independently). \\

  Operational complexity
    & One process, one port, one log stream.
    & Seven processes, seven ports, distributed tracing required for end-to-end request tracing. \\

  Delivery speed
    & Single-command start: \texttt{uvicorn app.main:app}.
    & Seven-command start abstracted by \texttt{run\_microservices.py} or Docker Compose. \\

  \bottomrule
\end{longtable}

\textbf{Trade-off summary:} The layered monolith is strictly superior for development velocity,
debugging simplicity, and in-process performance. The microservices decomposition becomes
valuable when individual services have genuinely different scaling needs or when teams grow large
enough that independent deployment cycles reduce coordination overhead. For the current team
size and stage, the monolith is the correct default; the microservices branch serves as a
forward migration target.

\subsection{NFR quantification}
\label{sec:nfr-quantification}

\begin{longtable}{p{2.6cm} p{3.0cm} p{3.0cm} p{7.4cm}}
  \toprule
  \textbf{NFR} & \textbf{Metric} & \textbf{Result} & \textbf{Method} \\
  \midrule
  \endhead

  Performance
    & Firestore reads per \texttt{GET /api/v1/meals} at \texttt{pageSize=100}
    & $\mathbf{101 \to 1}$ \textbf{reads}
    & Items are denormalized into the meal document at write time
      (meals service). At $N{=}100$: $101 \to 1$ ($100\times$ fewer reads). \\

  Performance (frontend bundle)
    & Initial JS bundle size (gzip)
    & \textbf{137.14 KB}
    & Route-level code splitting via \texttt{React.lazy()} and \texttt{<Suspense>} reduces the initial parse budget. \\

  Maintainability
    & Average cyclomatic complexity (CC) across \texttt{backend/app/}
    & \textbf{A (2.98)} average; single D (27) outlier
    & \texttt{radon cc backend/app -a --total-average};Sole outlier:
      \texttt{\_compute\_daily\_progress} in the analytics service (rank D, CC=27). Total SLOC: \textbf{2,393} across 52 files; largest file \texttt{recommendation\_engine.py} at 500 SLOC. \\

  Testability
    & Backend line coverage; test count; suite runtime
    & \textbf{76\% coverage}, \textbf{56/56 passing}, \textbf{\textasciitilde{}1.0 s}
    & \texttt{pytest --cov=app --cov-report=term-missing}: 56 tests across
      16 files, 1,809 statements covered, 430 missed. Suite runs in \textasciitilde{}1.0 s
      (mocked Firestore, no real I/O). CI enforces \texttt{--cov-fail-under=60}. \\

  Extensibility
    & Cost to add a new nutrition data source (OCP compliance)
    & \textbf{+1 class, 0 edits}
    & \texttt{test\_open\_closed\_extension\_via\_new\_provider} in the backend test suite
      plugs in a \texttt{SushiProvider} class without editing \texttt{IndianNutritionProvider},
      \texttt{GlobalCommonFoodsProvider}, \texttt{DefaultEstimateProvider}, or
      \texttt{NutritionResolver}. Test is green in CI; any future change that closes the
      chain will fail CI automatically. \\

  Observability
    & Prometheus metric families; log fields per request; request-ID coverage
    & \textbf{2 families}; \textbf{6 fields}; \textbf{100\%}
    & \texttt{/metrics} exposes: \texttt{http\_requests\_total} (counter) and
      \texttt{http\_request\_duration\_seconds} (histogram, 11 buckets: 5 ms--10 s).
      Every log line carries: \texttt{request\_id}, \texttt{method}, \texttt{path},
      \texttt{status\_code}, \texttt{duration\_ms}, \texttt{client\_ip}.
      Every request mints or inherits an \texttt{x-request-id}---100\% coverage
      by construction. In the \texttt{comparison} branch the gateway propagates the same ID
      to the upstream service, enabling cross-hop log joins. \\

  Scalability (routing layer)
    & Sustained RPS and p95 latency at the routing boundary (\texttt{/health}, no auth),
      monolith vs.\ gateway
    & Monolith: \textbf{73.8 RPS}, p95 \textbf{3.78 ms}; Gateway: \textbf{73.9 RPS},
      p95 \textbf{2.73 ms} (both at 0 errors)
    & \texttt{bench.py --rps 80 --duration 15} on both stacks; Both saturate the 80 RPS scheduler with zero errors.
      The gateway is marginally faster on \texttt{/health} because it carries less middleware
      (no rate-limiter, no domain-error handler). Qualitative scaling advantage of
      microservices---independent horizontal scaling per service---cannot be captured on a
      local single-machine bench. \\

  \bottomrule
\end{longtable}

% ===========================
% Reflection and Lessons Learned
% ===========================
\section{Reflection and Lessons Learned}

\begin{itemize}
  \item Layered architecture (routes / schemas / services / repositories) made testing fast and clean
  \item ADRs were useful for making and communicating decisions, and for future reference
  \item Nutrition accuracy remains unresolved
\end{itemize}

\end{document}
